\Needspace{9\baselineskip}
\originalchapter{级联码与阈值}
\originalsection{逻辑错误递推}
距离3容错块在局域独立小错误率$p$下, 单个故障可恢复, 失败至少需一对恶性故障. 若一个扩展门矩形有$A$对足以致错的位置, 粗略上界$p_1\le Ap^2$. 级联递推$p_{l+1}\le Ap_l^2$给$p_l\le(Ap)^{2^l}/A$. 当$p<1/A$时误差随层数双指数下降; 要总$L$个逻辑门失败小于$\epsilon$, 选$p_l\lesssim\epsilon/L$, 层数只随$\log\log(L/\epsilon)$增长, 空间时间开销成为对数的幂.

这不是所有设备共用的“阈值数值”. $A$依电路、辅助态验证、连接性和错误模型. 长程相关噪声可破坏独立两故障计数. 阈值定理的严格版本还控制坏矩形重叠和局域噪声范数; 上述推导是明确独立模型下的教学估计, 不冒充全部一般性证明.

\originalsection{综合征提取与辅助态}
若同一辅助比特顺序与许多数据耦合, 一个辅助X或Z错误可能传播成同块多错误. Shor式猫态把每个辅助分量只接一个数据位置, 可限制传播. 但猫态制备本身会产生相关错误, 必须验证或使用有相应故障保证的准备方法; “假定猫态最多一错误”属于题设, 不是实际制备结论. 测量综合征也可能错, 重复读数与时空解码可处理.

Bacon--Shor把码空间分为逻辑与规范自由度. 九比特3乘3格的局域双体规范算符可相乘得到六体稳定子, 因此无需直接耦合一个六体测量. 规范测量彼此可不对易, 但所提取稳定子为中心, 测量次序与多轮推断仍要设计. 规范自由度的错误可忽略, 逻辑错误不可; 这个区别使某些多位置错误等价于低权重规范类, 不表示任意两物理错误都可纠正.
